\section{拉回是双射}\label{sec:pullback-is-bijection}

若 $X\subseteq\mathbb{A}^{m}$ 和 $Y\subseteq\mathbb{A}^{n}$ 为仿射簇, 则把多项式映射送到其拉回的对应给出了两个态射集合之间的函数:
\[
  \begin{aligned}
    \{X\to Y\text{ 的多项式映射}\}
      &\to\{K[Y]\to K[X]\text{ 的 }K\text{-代数同态}\}\\
    F&\mapsto F^{*}.
  \end{aligned}
\]
本节的主要结论是这个函数为双射.

\begin{Proposition}{多项式函数的拉回给出双射}\leavevmode\label{prop:pullback-bijection}
  设 $X\subseteq\mathbb{A}^{m}$ 和 $Y\subseteq\mathbb{A}^{n}$ 为仿射簇. 对应 $F\mapsto F^{*}$ 是从多项式映射 $X\to Y$ 的集合到 $K$-代数同态 $K[Y]\to K[X]$ 的集合的双射.
\end{Proposition}

\BookPairedNoteLeft
  {用范畴论的语言来说, Proposition~\ref{prop:pullback-bijection} 所给的双射表明, 把 $X$ 送到 $K[X]$, 把 $F$ 送到 $F^{*}$ 的函子是\textquotedbl{}全忠实的\textquotedbl{}.}
  {为证明 Proposition~\ref{prop:pullback-bijection}, 只需构造把 $F$ 送到 $F^{*}$ 的过程的逆. 也就是说, 给定任意 $K$-代数同态 $\phi:K[Y]\to K[X]$, 只需说明存在唯一的多项式映射 $F:X\to Y$ 使 $F^{*}=\phi$.}
\noindent 这一结论的证明由于记号的缘故可能不易读懂, 因此先看一个具体例子. 建议读者在学习证明时回看这一例子.

\begin{Example}{由拉回恢复原映射}\leavevmode\label{ex:invert-pullback}
  考虑 Example~\ref{ex:parabola-map-to-hyperbolic-paraboloid} 中的仿射簇:
  \[
    X=\mathcal{V}(x_{2}-x_{1}^{2})\subseteq\mathbb{A}^{2}
    \qquad\text{且}\qquad
    Y=\mathcal{V}(y_{1}y_{2}-y_{3})\subseteq\mathbb{A}^{3}.
  \]
  利用 $x_{2}-x_{1}^{2}$ 和 $y_{1}y_{2}-y_{3}$ 都不可约这一事实, 得到
  \[
    K[X]=\frac{K[x_{1},x_{2}]}{\langle x_{2}-x_{1}^{2}\rangle}
    \qquad\text{且}\qquad
    K[Y]=\frac{K[y_{1},y_{2},y_{3}]}{\langle y_{1}y_{2}-y_{3}\rangle}.
  \]
  考虑 $K$-代数同态 $\phi:K[Y]\to K[X]$. 这样的同态由每个生成元 $[y_{i}]$ 的像 $[f_{i}]$ 决定, 其中 $f_{i}\in K[x_{1},x_{2}]$ 是某个多项式. 例如, 考虑同态
  \[
    \begin{aligned}
      \phi:K[Y]&\to K[X]\\
      \phi([y_{1}])&=[x_{1}+x_{2}]\\
      \phi([y_{2}])&=[x_{1}]\\
      \phi([y_{3}])&=[x_{1}x_{2}+x_{2}].
    \end{aligned}
  \]
  $\phi$ 是 $K$-代数同态, 这一事实决定了任意元素 $[g]\in K[Y]$ 的像. 例如,
  \[
    \phi([y_{1}^{2}+y_{2}-y_{2}y_{3}])=[(x_{1}+x_{2})^{2}+x_{1}-x_{1}(x_{1}x_{2}+x_{2})].
  \]
  更一般地, 对任意 $g\in K[y_{1},y_{2},y_{3}]$, 有
  \[
    \phi([g])=[g(x_{1}+x_{2},x_{1},x_{1}x_{2}+x_{2})].
  \]
  并不是任取三个多项式 $f_{1},f_{2},f_{3}$ 都能得到良定义的 $K$-代数同态. 必须有 $[0]=[y_{1}y_{2}-y_{3}]$ 被送到 $[0]\in K[X]$, 或者说 $f_{1}f_{2}-f_{3}$ 必须属于 $\langle x_{2}-x_{1}^{2}\rangle$. 我们所选的 $f_{1},f_{2},f_{3}$ 的确满足这一点:
  \begin{equation}\label{eq:inverse-pullback-relation-in-ideal}
    (x_{1}+x_{2})x_{1}-(x_{1}x_{2}+x_{2})=x_{1}^{2}-x_{2}\in\langle x_{2}-x_{1}^{2}\rangle.
  \end{equation}

  给出 $\phi$ 之后, 能否找到多项式映射 $F:X\to Y$, 使 $F^{*}=\phi$? 这样的映射把 $X\subseteq\mathbb{A}^{2}$ 的元素送到 $Y\subseteq\mathbb{A}^{3}$ 的元素, 因此它由三个二元多项式定义. 能想到哪三个多项式呢? 有一种自然的选择: 定义 $\phi$ 时所用的三个多项式 $f_{1},f_{2},f_{3}$. 换句话说, 考虑函数
  \[
    \begin{aligned}
      F:X&\to Y\\
      F(a_{1},a_{2})&=(a_{1}+a_{2},a_{1},a_{1}a_{2}+a_{2}).
    \end{aligned}
  \]
  虽然 $F$ 显然是多项式映射, 但仍应确认它确实把 $X$ 的元素送到 $Y$ 的元素, 即 $F(a_{1},a_{2})$ 满足 $Y$ 的定义方程. 具体地, 验证
  \begin{equation}\label{eq:inverse-pullback-map-satisfies-relation}
    (a_{1}+a_{2})a_{1}-(a_{1}a_{2}+a_{2})=a_{1}^{2}-a_{2}=0,
  \end{equation}
  其中最后一个等号成立是因为 $(a_{1},a_{2})\in X=\mathcal{V}(x_{2}-x_{1}^{2})$. 注意 \eqref{eq:inverse-pullback-relation-in-ideal} 与 \eqref{eq:inverse-pullback-map-satisfies-relation} 的相似之处: 验证 $\phi$ 良定义所需的条件, 恰好就是验证 $F$ 将 $X$ 映入 $Y$ 所需的条件.

  最后, 为确认 $F^{*}=\phi$, 只需检查这两个同态在 $K[Y]$ 的生成元 $[y_{1}],[y_{2}],[y_{3}]$ 上相同. 把 $[y_{i}]$ 看成函数 $Y\to K$, 它就是坐标函数
  \[
    \begin{aligned}
      [y_{i}]:Y&\to K\\
      [y_{i}](b_{1},b_{2},b_{3})&=b_{i}.
    \end{aligned}
  \]
  因此, 按拉回的定义,
  \[
    \begin{aligned}
      F^{*}([y_{1}])(a_{1},a_{2})
        &=([y_{1}]\circ F)(a_{1},a_{2})\\
        &=[y_{1}](a_{1}+a_{2},a_{1},a_{1}a_{2}+a_{2})=a_{1}+a_{2}.
    \end{aligned}
  \]
  换句话说,
  \[
    F^{*}([y_{1}])=[x_{1}+x_{2}],
  \]
  从而 $F^{*}([y_{1}])=\phi([y_{1}])$. 类似地, $F^{*}$ 和 $\phi$ 在另外两个生成元上也相同, 所以我们已经构造出一个多项式映射 $F$, 使 $F^{*}=\phi$. 此外, 这一构造实质上也说明了 $F$ 的唯一性: 要求 $F^{*}([y_{i}])=\phi([y_{i}])$ 决定了 $F$ 的第 $i$ 个分量函数, 而这些分量函数唯一决定 $F$.
\end{Example}

推广这一例子的思路, 现在可以证明 Proposition~\ref{prop:pullback-bijection}.

\begin{Proof}\textbf{of Proposition~\ref{prop:pullback-bijection}.}\quad
  设 $X\subseteq\mathbb{A}^{m}$ 和 $Y\subseteq\mathbb{A}^{n}$ 为仿射簇, 有
  \[
    K[X]=\frac{K[x_{1},\ldots,x_{m}]}{\mathcal{I}(X)}
    \qquad\text{且}\qquad
    K[Y]=\frac{K[y_{1},\ldots,y_{n}]}{\mathcal{I}(Y)}.
  \]
  为证明该命题, 我们说明对于每个 $K$-代数同态 $\phi:K[Y]\to K[X]$, 存在唯一的多项式映射 $F:X\to Y$ 使 $F^{*}=\phi$.

  (存在性.) 设 $\phi:K[Y]\to K[X]$ 为 $K$-代数同态. 沿用上例的构造, 令 $F_{i}=\phi([y_{i}])\in K[X]$, 考虑多项式映射
  \[
    F=(F_{1},\ldots,F_{n}):X\to\mathbb{A}^{n}.
  \]
  与上例一样, 必须验证 $F$ 的像落在 $Y$ 中, 从而得到多项式映射 $F:X\to Y$. 为此, 设 $a\in X$; 必须证明
  \[
    (F_{1}(a),\ldots,F_{n}(a))\in Y.
  \]
  由于 $Y=\mathcal{V}(\mathcal{I}(Y))$, 只需验证对所有 $h\in\mathcal{I}(Y)$, 有
  \[
    h(F_{1}(a),\ldots,F_{n}(a))=0.
  \]
  设 $h\in\mathcal{I}(Y)$. 则 $[h]=0\in K[Y]$, 因为 $\phi$ 是同态, 所以 $\phi([h])=0\in K[X]$. 但这蕴含
  \[
    \phi([h(y_{1},\ldots,y_{n})])=h(\phi([y_{1}]),\ldots,\phi([y_{n}]))=h(F_{1},\ldots,F_{n})
  \]
  是 $X$ 上的零函数, 从而 $h(F_{1}(a),\ldots,F_{n}(a))=0$. 因此 $F$ 确实是从 $X$ 到 $Y$ 的多项式映射.

  还需证明 $F^{*}=\phi$. 与上例一样, 对每个 $i=1,\ldots,n$,
  \[
    F^{*}([y_{i}])=[y_{i}]\circ(F_{1},\ldots,F_{n})=F_{i}=\phi([y_{i}]).
  \]
  于是 $F^{*}$ 和 $\phi$ 在 $K[Y]$ 的生成元 $[y_{i}]$ 上相同, 从而 $F^{*}=\phi$.

  (唯一性.) 设 $F,G:X\to Y$ 为满足 $F^{*}=G^{*}$ 的多项式映射. 必须证明 $F=G$. 按定义,
  \[
    F=(F_{1},\ldots,F_{n})\qquad\text{且}\qquad G=(G_{1},\ldots,G_{n}),
  \]
  其中 $F_{i},G_{i}\in K[X]$ 是 $X$ 上的多项式函数. 在 $[y_{i}]$ 处对 $F^{*}$ 和 $G^{*}$ 求值, 有
  \[
    F^{*}([y_{i}])=[y_{i}]\circ(F_{1},\ldots,F_{n})=F_{i}\in K[X],
  \]
  类似地, $G^{*}([y_{i}])=G_{i}\in K[X]$. 因为 $F^{*}=G^{*}$, 所以
  \[
    F_{i}=F^{*}([y_{i}])=G^{*}([y_{i}])=G_{i}
  \]
  对所有 $i=1,\ldots,n$ 成立. 因此 $F=G$.
\end{Proof}

证明 Proposition~\ref{prop:pullback-bijection} 所作的工作给我们带来的结果是: 现在可以通过研究两个仿射簇的坐标环, 准确判定它们是否同构.

\begin{Corollary}{坐标环判定同构}\leavevmode\label{cor:coordinate-rings-detect-isomorphisms}
  设 $X\subseteq\mathbb{A}^{m}$ 和 $Y\subseteq\mathbb{A}^{n}$ 为仿射簇. 则 $F:X\to Y$ 是同构当且仅当 $F^{*}:K[Y]\to K[X]$ 是同构, 并且
  \[
    X\cong Y\qquad\Longleftrightarrow\qquad K[X]\cong K[Y].
  \]
\end{Corollary}

\begin{Proof}
  第一个断言的正向蕴含正是 Corollary~\ref{cor:pullback-of-isomorphism} 的内容. 为证明逆向蕴含, 假设 $F^{*}:K[Y]\to K[X]$ 是同构, 其逆为 $(F^{*})^{-1}:K[X]\to K[Y]$. 由 Theorem~\ref{prop:pullback-bijection} 中的满射性, 存在多项式映射 $G:Y\to X$, 使 $G^{*}=(F^{*})^{-1}$. 我们断言 $G=F^{-1}$. 的确, 根据 Proposition~\ref{prop:pullback-preserves-compositions-and-identity},
  \[
    (F\circ G)^{*}=G^{*}\circ F^{*}=(F^{*})^{-1}\circ F^{*}
      =\operatorname{id}_{K[Y]}=(\operatorname{id}_{Y})^{*}.
  \]
  因为 $F\circ G$ 和 $\operatorname{id}_{Y}$ 有相同的拉回, Theorem~\ref{prop:pullback-bijection} 中的单射性蕴含 $F\circ G=\operatorname{id}_{Y}$. 类似地, $G\circ F=\operatorname{id}_{X}$. 所以 $G=F^{-1}$, 从而 $F$ 是同构.

  现在考虑第二个断言. 正向蕴含是第一个断言的推论: 若 $X\cong Y$, 则存在同构 $F:X\to Y$, 其拉回给出同构 $F^{*}:K[Y]\to K[X]$, 这说明 $K[X]\cong K[Y]$. 逆向蕴含也来自第一个断言, 只须利用 Theorem~\ref{prop:pullback-bijection} 中的满射性, 注意 $K[X]$ 与 $K[Y]$ 之间的任意同构必定是 $X$ 与 $Y$ 之间某个多项式映射的拉回.
\end{Proof}

\BookExerciseHeading
\begin{BookExercises}
  \item\label{ex:inverse-pullback-of-natural-inclusion} 假设 $m\leq n$, 考虑自然单射
    \[
      \phi:K[x_{1},\ldots,x_{m}]\to K[x_{1},\ldots,x_{n}].
    \]
    描述相应的多项式映射 $F:\mathbb{A}^{n}\to\mathbb{A}^{m}$.
  \item\label{ex:inverse-pullback-of-specialization} 假设 $m\leq n$, 考虑满射
    \[
      \begin{aligned}
        \phi:K[x_{1},\ldots,x_{n}]&\to K[x_{1},\ldots,x_{m}]\\
        f(x_{1},\ldots,x_{n})&\mapsto f(x_{1},\ldots,x_{m},0,\ldots,0).
      \end{aligned}
    \]
    描述相应的多项式映射 $F:\mathbb{A}^{m}\to\mathbb{A}^{n}$.
  \item\label{ex:inverse-pullback-into-affine-plane} 考虑 $K$-代数同态
    \[
      \phi:K[x,y]\to K[t],
    \]
    其定义为
    \[
      \phi(f)=f(t+1,t^{2}+t).
    \]
    对于哪些仿射簇 $X,Y$ 以及哪个多项式映射 $F:X\to Y$, 有 $\phi=F^{*}$?
  \item\label{ex:inverse-pullback-from-quadric} 设 $X=\mathcal{V}(x^{2}+y^{2}+z^{2}-1)\subseteq\mathbb{A}^{3}_{\mathbb{C}}$, 有
    \[
      \mathbb{C}[X]=\frac{\mathbb{C}[x,y,z]}{\langle x^{2}+y^{2}+z^{2}-1\rangle}.
    \]
    考虑 $\mathbb{C}$-代数同态
    \[
      \begin{aligned}
        \phi:\mathbb{C}[u,v]&\to\mathbb{C}[X]\\
        \phi(f)&=[f(x+y+z,xyz)].
      \end{aligned}
    \]
    对于哪个多项式映射 $F:X\to\mathbb{A}^{2}$, 有 $\phi=F^{*}$?
  \item\label{ex:pullback-bijection-for-plane} 设 $f_{1},f_{2},f_{3}\in K[w]$, 且 $X=\mathcal{V}(x+y-z)\subseteq\mathbb{A}^{3}$.
    \begin{enumerate}[label=(\alph*)]
      \item $f_{1},f_{2},f_{3}$ 满足什么条件时,
        \[
          \begin{aligned}
            F:\mathbb{A}^{1}&\to X\\
            F(a)&=(f_{1}(a),f_{2}(a),f_{3}(a))
          \end{aligned}
        \]
        给出一个到 $X$ 的良定义多项式映射? 给出满足和不满足该条件的 $f_{1},f_{2},f_{3}\in K[w]$ 各一个具体例子.
      \item $f_{1},f_{2},f_{3}$ 满足什么条件时, 存在良定义的 $K$-代数同态
        \[
          \phi:\frac{K[x,y,z]}{\langle x+y-z\rangle}\to K[w],
        \]
        它在生成元上的定义为
        \[
          \begin{aligned}
            \phi([x])&=f_{1}(w)\\
            \phi([y])&=f_{2}(w)\\
            \phi([z])&=f_{3}(w)?
          \end{aligned}
        \]
        给出满足和不满足该条件的 $f_{1},f_{2},f_{3}\in K[w]$ 各一个具体例子.
      \item 利用 (a) 和 (b), 具体描述多项式映射 $\mathbb{A}^{1}\to X$ 的集合与 $K$-代数同态 $K[X]\to K[\mathbb{A}^{1}]$ 的集合之间的双射.
    \end{enumerate}
  \item\label{ex:pullback-to-a-point} 设 $X\subseteq\mathbb{A}^{m}$ 为仿射簇, $Y\subseteq\mathbb{A}^{n}$ 为单点. 从 $X$ 到 $Y$ 只有一个可能的多项式映射, 因而由 Theorem~\ref{prop:pullback-bijection}, 只有一个可能的 $K$-代数同态 $K[Y]\to K[X]$. $K[Y]$ 是什么? 这个唯一的 $K$-代数同态 $K[Y]\to K[X]$ 是什么?
  \item\label{ex:pullback-from-irreducible-to-two-points} 假设 $X\subseteq\mathbb{A}^{m}$ 为不可约仿射簇, $Y\subseteq\mathbb{A}^{n}$ 由两个不同的点组成. 证明恰有两个多项式映射 $F:X\to Y$. 具体描述这两个映射及其拉回.
  \item\label{ex:pullbacks-of-polynomial-functions} 若 $X\subseteq\mathbb{A}^{n}$ 是任意仿射簇, 则多项式映射 $X\to\mathbb{A}^{1}$ 的集合恰是 $K[X]$, 因而由 Theorem~\ref{prop:pullback-bijection}, $K[X]$ 与 $K$-代数同态 $K[\mathbb{A}^{1}]\to K[X]$ 的集合之间有一个双射. 具体描述这个双射.
\end{BookExercises}
